% =====================================================================
% QUESTION 4 — File I/O  (5 marks)
% Demonstrates: plain prose + task bullets + before/after file table
% =====================================================================
\question{4}{5}

\noindent The text file \texttt{src/transactions.txt} contains several lines,
one purchase per line. Each line stores a customer ID (string), the customer
name (string), and the purchase amount (double), separated by a single space.
A sample input is shown in the left column below.

\noindent Write a Java program that:
\begin{itemize}[itemsep=2pt, topsep=4pt]
    \item reads \texttt{src/transactions.txt} line by line, \hfill [1]
    \item writes every customer ID and name to \texttt{src/customers.txt}, \hfill [1]
    \item writes the customer name and amount to \texttt{src/summary.txt},
          but only for purchases above \texttt{5000.0}, \hfill [2]
    \item prints ``Processed N records.'' once all records are read. \hfill [1]
\end{itemize}

\noindent For the sample input below, your program must produce the output
file shown on the right (purchases that are skipped appear as ---).

\vspace{0.3cm}
% plain tabular (not tabularx): this table has no X column, and tabularx's
% natural-width measuring pass overflows even when the columns fit.
{\renewcommand{\arraystretch}{1.35}
\begin{tabular}{|>{\raggedright\arraybackslash\ttfamily\small}p{0.40\textwidth}|>{\raggedright\arraybackslash\ttfamily\small}p{0.50\textwidth}|}
\hline
\textbf{\texttt{transactions.txt} (input)} & \textbf{\texttt{summary.txt} (output)} \\ \hline
C-101 Ayesha 2400.0  & --- \\ \hline
C-102 Tanvir 12400.0 & Tanvir 12400.0 \\ \hline
C-103 Rakib 8700.0   & Rakib 8700.0 \\ \hline
C-104 Samia 15600.0  & Samia 15600.0 \\ \hline
C-105 Nadia 3400.0   & --- \\ \hline
C-106 Farhan 1900.0  & --- \\ \hline
\end{tabular}}
